\medskip

\textbf{Partie A}

\medskip

\begin{enumerate}
\item 
\begin{itemize}
\item Le point M est défini par $\vectt{BM} = \vectt{AB}$, donc il appartient à la droite (AB), et donc au plan (FAB). On en déduit que la droite (FM) est une droite incluse dans le plan (FAB).
\item Le plan (FAB) contient la face ABFE du cube. Comme les faces du cube sont des carrés, on en déduit que la droite (FG) est perpendiculaire  à la droite (EF) (via le carré EFGH) et à la droite (FB) (via le carré BCGF).

Ainsi, (FG) est orthogonale à deux droites sécantes du plan (FAB) : on en déduit que (FG) est orthogonale au plan (FAB).
\item Puisque (FG) est orthogonale à (FAB), elle est donc également orthogonale à toute droite du plan (FAB), notamment la droite (FM).
\item Comme F est un point d'intersection entre (FB) et (FM), les droites sont orthogonales et sécantes, on peut donc conclure que les droites sont bien perpendiculaires.
		\end{itemize}

\item 
\begin{itemize}
\item On a déjà établi que M est sur la droite (AB), avec : $\vectt{AM} = \vectt{AB} + \vectt{BM} = 2\vectt{AB}$;
\item Comme ABCDEFGH est un cube, on a notamment : $\vectt{HG} = \vectt{AB}$;
\item On en déduit donc : \quad $\vectt{AM} = 2\vectt{HG}$;
\item Les droites (AM) et (HG) sont dirigées par des vecteurs directeurs colinéaires, donc elles sont parallèles, et donc coplanaires;
\item cela implique que les quatre points A, M, H et G sont également coplanaires.
\end{itemize}
\end{enumerate}

\medskip

\textbf{Partie B}

\medskip

Le repère $\left(\mathrm{A} ~;~ \vectt{AB}, \vectt{AD}, \vectt{AE}\right)$ est un repère orthonormé de l'espace : les vecteurs de base sont deux à deux orthogonaux et de même norme, car ils sont basés sur les arêtes consécutives d'un cube.

On a alors les coordonnées : A$(0~;~0~;~0)$; \quad  B$(1~;~0~;~0)$; \quad  C$(1~;~1~;~0)$; \quad  D$(0~;~1~;~0)$; \quad  E$(0~;~0~;~1)$; \quad  F$(1~;~0~;~1)$; \quad  G$(1~;~1~;~1)$; \quad  H$(0~;~1~;~1)$ ~~et ~~ M$(2~;~0~;~0)$.

\begin{enumerate}
		\item On a :\quad $\vectt{GM}~\begin{pmatrix}2 - 1\\0-1\\0-1\end{pmatrix}= \begin{pmatrix}1\\-1\\-1\end{pmatrix}$

		et, comme A est l'origine du repère, les coordonnées de H sont aussi celles du vecteur $\vectt{AH} ~\begin{pmatrix}0\\1\\1\end{pmatrix}$.

		Ces coordonnées sont clairement non proportionnelles, donc les vecteurs ne sont pas colinéaires.

		\item
		\begin{enumerate}
			\item (GM) passe par le point G de coordonnées $(1~;~1~;~1)$ et est dirigée par $\vectt{GM}\begin{pmatrix}1\\-1\\-1\end{pmatrix}$.

			Une représentation paramétrique de (GM) est donc :

			$\left\{\begin{array}{l}x = x_{\mathrm{G}} + tx_{\vectt{GM}}\\y = y_{\mathrm{G}} + ty_{\vectt{GM}}\\z = z_{\mathrm{G}} + tz_{\vectt{GM}}\end{array}\right. \quad \text{avec } t \in \R$ \qquad soit \qquad
			$\left\{\begin{array}{l}x = 1 + t\\y = 1 - t\\z = 1 - t\end{array}\right. \quad \text{avec } t \in \R$

			\item Cherchons s'il est possible qu'un point de paramètre $t$ sur (GM) ait les mêmes coordonnées qu'un point de paramètre $k$ sur (AH). Pour cela, résolvons le système :

			$S : \left\{\begin{array}{l}1 + t = 0\\1 - t = k\\1 - t = k\end{array}\right.$
			\qquad on a :\quad
			$S \iff \left\{\begin{array}{l}t = -1 \\k = 1 - (-1)\\k = 1 - (-1)\end{array}\right.$
			\qquad
			Soit \quad$S \iff \left\{\begin{array}{l}t = -1 \\k = 2\\k = 2\end{array}\right.$

			On en déduit que seul un point est commun aux deux droites, c'est le point de paramètre $t = -1$ sur (GM), qui est aussi le point de paramètre $k = 2$ sur (AH).

			Qu'on utilise une représentation paramétrique ou l'autre pour trouver les coordonnées de N, on a bien :\quad N$\big(1 + (- 1) ~;~ 1 -(- 1) ~;~ 1 -(- 1)\big)$ \quad c'est-à-dire \quad N$(0 ~;~ 2 ~;~ 2)$.
		\end{enumerate}

		\item
		\begin{enumerate}
			\item On a :\quad $\vectt{AM} = \begin{pmatrix}2\\0\\0\end{pmatrix}$,\quad et \quad $\vectt{AN} = \begin{pmatrix}0\\2\\2\end{pmatrix}$.

			Comme le repère est orthonormal, on peut calculer le produit scalaire à l'aide des coordonnées :
			\quad $\vectt{AM} \cdot \vectt{AN} = 2 \times 0 + 0 \times 2 + 0 \times 2 = 0$.

			Les vecteurs sont donc orthogonaux, donc les droites (AM) et (AN) sont orthogonales (et sécantes en A), donc perpendiculaires :

			Le triangle AMN est bien rectangle en A.

			\item Le repère est orthonormal donc on peut aussi calculer des longueurs avec les coordonnées :

			$\mathrm{AM} = \sqrt{2^2 + 0^2 + 0^2} = \sqrt{4} = 2$,\quad $\mathrm{AN} = \sqrt{0^2 + 2^2 + 2^2} = \sqrt{8} = 2\sqrt{2}$.

			Avec [AM] pour base, la hauteur correspondante est [AN] puisque le triangle est rectangle en A.

			On a donc :\quad $\mathcal{A}_{\mathrm{AMN}}= \dfrac{\mathrm{AM} \times \mathrm{AN}}{2} = \dfrac{2\times 2\sqrt{2}}{2} = 2\sqrt{2}$.
		\end{enumerate}

		\item
		\begin{enumerate}
			\item J est le centre de BCGF, donc notamment le milieu de [BG] :

			$\mathrm{J} = \left(\dfrac{x_\mathrm{B} + x_{\mathrm{G}}}{2} ~;~ \dfrac{y_\mathrm{B} + y_{\mathrm{G}}}{2} ~;~ \dfrac{z_\mathrm{B} + z_{\mathrm{G}}}{2}\right)$,
			\quad soit \quad J$\left(\dfrac{1 + 1}{2} ~;~ \dfrac{0 + 1}{2}~;~ \dfrac{0 + 1}{2}\right)$.

			Finalement, les coordonnées de J sont : \quad J$(1 ~;~ 0,5 ~;~ 0,5)$.

			\item On a :\quad $\vectt{FJ} \begin{pmatrix} 1 - 1\\0,5 - 0\\0,5 - 1 \end{pmatrix}$\quad soit :\quad $\vectt{FJ} \begin{pmatrix} 0\\0,5\\-0,5\end{pmatrix}$.

			Vérifions : \quad $\vectt{FJ} \cdot \vectt{AM} = 0 \times 2 + 0,5 \times 0 + (-0,5) \times 0 = 0$.

			et, par ailleurs :\quad $\vectt{FJ} \cdot \vectt{AN} = 0 \times 0 + 0,5 \times 2 + (-0,5) \times 2 = 0$.

			Donc $\vectt{FJ}$ est orthogonal à deux vecteurs formant une base de (AMN) : c'est bien un vecteur normal au plan (AMN).

			\item Puisque $\vectt{FJ}$ est normal à (AMN), on en déduit qu'(AMN) admet une équation de la forme :\quad $0\times x + 1\times y + (- 1)\times z + d = 0$,~~ soit ~~$y - z + d = 0$ \quad où $d \in \R$.

			Comme A$(0 ~;~ 0 ~;~ 0)$ est sur (AMN), ses coordonnées doivent vérifier l'équation, donc :

			$0 - 0 + d = 0 \iff d = 0$.

			Finalement, une équation de (AMN) est $y - z = 0$.

			En testant l'appartenance de J à (AMN) :\quad $\begin{aligned}[t] y_{\mathrm{J}} - z_{\mathrm{J}} &= 0,5 - 0,5\\ &= 0
			\end{aligned}$

			Les coordonnées de J vérifient l'équation du plan, donc J appartient au plan (AMN).

			Comme $\vectt{FJ}$ est orthogonal au plan et que J est dans le plan, J est bien le projeté orthogonal de F sur (AMN).
		\end{enumerate}

		\item Pour déterminer le volume du tétraèdre AMNF, on choisir comme base le triangle AMN, la hauteur correspondante étant [FJ].

		$\mathcal{V}_{\mathrm{AMNF}} = \dfrac{1}{3} \times 2\sqrt{2} \times \mathrm{FJ}$ \quad où \quad $\mathrm{FJ} = \sqrt{0^2 + 0,5^2 + 0,5^2} = \sqrt{0,5} = \dfrac{\sqrt{2}}{2}$.

		On a donc :\quad $\mathcal{V}_{\mathrm{AMNF}} = \dfrac{1}{3} \times 2\sqrt{2} \times \dfrac{\sqrt{2}}{2} = \dfrac{1}{3} \times 2 = \dfrac{2}{3}$.

		Pour la pyramide BCGFM, la base (seule face non triangulaire) le carré BCGF (carré de côté 1, donc d'aire 1), la hauteur étant la distance entre M et le plan (BCG).

		Comme on a déjà établi que M est un point de (AB), et que ABCDEFGH est un cube, on en déduit que (BM) est perpendiculaire au plan (BCG) et que B est donc le projeté orthogonal de M sur (BCG). La distance entre M et (BCG) est donc BM = AB = 1.

		$\mathcal{V}_{\mathrm{BCGFM}} = \dfrac{1}{3} \times 1 \times 1 = \dfrac{1}{3}$.

		Donc on a bien :\quad $\mathcal{V}_{\mathrm{AMNF}} = 2 \times \mathcal{V}_{\mathrm{BCGFM}}$.
	\end{enumerate}

	\bigskip
